\documentclass{article}
\usepackage[utf8]{inputenc}
\usepackage[margin=0.75in]{geometry}
\usepackage[dvipsnames]{xcolor} % adding colour
\usepackage{amsmath} % adding math equations
\usepackage{natbib} % adding BibTeX support
\usepackage{hyperref}



\begin{document}
	\begin{center}
    
    	% MAKE SURE YOU TAKE OUT THE SQUARE BRACKETS
    
		\LARGE{\textbf{STA380 Project Proposal}} \\
        \vspace{1em}
        \Large{Simulation-Based Visualization of Regression and Time Series Model Behavior} \\
        \vspace{1em}
        \normalsize\textbf{Names: Kevaan Buch, Kevin Joseph, Eshan Hussain, Zhaoqi Wang, Amin Amir} \\
        \normalsize{Emails: eshan.hussain@utoronto.ca, kevinj.joseph@mail.utoronto.ca, kevaan.buch@mail.utoronto.ca, lieyanxuegao.wang@mail.utoronto.ca, amin.amir@mail.utoronto.ca} \\
        \vspace{1em}
        \normalsize{University of Toronto, Mississauga}
     
	\end{center}
    \begin{normalsize}
    
    	\section{Project Topic}
        Simulation-Based Visualization of Regression and Time Series Model Behavior.

        \vspace{1em}

        \noindent{The objective of this project is to develop an interactive simulation and visualization tool that helps users build intuition for regression and time series models by exploring how they behave under correct and incorrect modelling assumptions. The project focuses on models with a single scalar-valued response variable to enable clear graphical visualization. We consider simple linear regression, autoregressive (AR), and moving average (MA) models, and use simulation and Monte Carlo repetition to study underfitting, residual dependence, and model mismatch.}
   


        \section{Simulation vs. Dataset}
        The results will be from using pure simulations	
        
        \section{Project Details}
        This project comprises a Shiny-based application with two primary components: single-run visualization and Monte Carlo summarization.
        \vspace{1em}

        \vspace{1em}
        \noindent{\textbf{Data Generation}}:
        All stochastic components will be generated using Gaussian white noise (i.i.d. Normal errors with mean zero and constant variance) through the acceptance-rejection algorithm.
        
        \vspace{1em}
        
        \noindent{Each data-generating process will be implemented directly from its defining equations. Linear regression processes will be generated by specifying an underlying relationship and adding simulated random noise. In contrast, AR and MA processes will be generated by iteratively updating values over time using past shocks.}
        \vspace{1em}
        \vspace{1em}

        \noindent{\textbf{Model Fitting}}: The intention is to allow users to explore both correctly specified and misspecified modeling scenarios.
        \vspace{1em}

        \noindent{Model fitting will be implemented explicitly by coding the underlying estimation procedures ourselves, rather than relying primarily on built-in R functions. This is intended to reinforce our understanding of how parameter estimates and residuals arise under each model.}

        
        \section{User Inputs (Shiny Components)}
           The user will be presented with a menu where they can select the data-generating process and the model family, as well as any additional parameters they can use. Following this, they will have a visualization menu, where they can toggle parameters about the graph itself.

        \vspace{1em}
        \noindent{Users will be able to set the following:}

        \begin{enumerate}
            \item \textbf{General Data Generating Process (DGP) Input Controls:} 
                \begin{itemize}
                    \item Sample size (n)
                    \item Noise variance
                    \item  Random seed (optional)
                \end{itemize}
                
             \item \textbf{Linear Regression DGP Input Controls}: 
                \begin{itemize}
                    \item Coefficient values (optional)
                    \item Number of coefficients (optional)
                    \item  Functional form of the DGP (optional)
                \end{itemize}
            \item \textbf{MA(n) DGP Input Controls}: 
                \begin{itemize}
                    \item Order of the MA process
                    \item coefficient values (optional)
                \end{itemize}
            \item \textbf{AR(n) DGP Input Controls}: 
                \begin{itemize}
                    \item Order of the AR process
                    \item Coefficient values (optional)
                \end{itemize}
            \item \textbf{Model Fitting}: 
                \begin{itemize}
                    \item Users will be able to select the model family (linear regression, AR(n), MA(n))
                    \item If an AR(n) or a MA(n)) process was selected, the user will specify the order of the respective process
                \end{itemize}
            \item \textbf{Visualization Toggles}: 
                \begin{itemize}
                    \item Data and fitted values
                    \item Residual time plot
                    \item Autocorrelation function of the residuals
                    \item Histogram of the residuals
                    \item  Change/specify colours used
                \end{itemize}
        \end{enumerate}

\end{normalsize}

\nocite{*}
\bibliographystyle{abbrvnat}
\bibliography{bibliography}
\vspace{0.8cm}

\end{document}